\section{Bayesian updating and prediction: complete proofs}
\label{sec:annex-p15-bayesian-updating-prediction}

Use the parameter, observation and kernel conventions of
\cref{sec:p15-bayesian-updating-prediction}. All integrations of
nonnegative functions below are justified before signed differences
are taken. Probability under the joint prior model is distinguished
from probability under a fixed-parameter sampling law.

\begin{annexitem}{annex:p15.bayes-posterior-kernel}{thm:p15-bayes-posterior-kernel}{The joint measure, exceptional observations and posterior integrals}
\begin{proof}
The product measure \(\Pi\otimes\nu\) exists and is sigma-finite.
Define \(Q\) on the product sigma-field by its density
\(f_\theta(x)\). Tonelli's theorem gives
\[
 \int f_\theta(x)\,(\Pi\otimes\nu)(d\theta,dx)
   =\int_\Theta\left(\int_{\mathcal X}f_\theta(x)\,\nu(dx)\right)
                                                    \Pi(d\theta)=1.
\]
This proves that \(Q\) is a probability measure with the stated
rectangle probabilities. Integrating in the other order proves
measurability of \(m\), \(\int m\,d\nu=1\), and the claimed marginal.
The set \(\{m=\infty\}\) has \(\nu\)-measure zero; otherwise the
integral of \(m\) would be infinite. Its integral is therefore zero.
The integral over \(\{m=0\}\) is zero as well. Consequently
the marginal measure \(\mu=m\nu\) satisfies \(\mu(G)=1\).

For fixed \(x\in G\), the numerator
\(A\mapsto\int_A f_\theta(x)\Pi(d\theta)\) is a finite
measure of total mass \(m(x)\). Division gives a probability measure.
Off \(G\), \(\Pi_x=\Pi\) is already a probability measure.
For each measurable \(A\), Tonelli proves that the numerator is
a measurable function of \(x\); division on \(G\) and the fixed
value off \(G\) preserve measurability. Thus \(\Pi_x\) is a kernel.

On \(G\), \(\Pi_x(A)m(x)=\int_A f_\theta(x)\Pi(d\theta)\).
The integral of the right side over \(G^c\) vanishes, since it is
bounded by the marginal integral there. Integrating over \(B\)
proves the rectangle identity. To pass from rectangles to the full
product sigma-field, let \(\mathcal E\) be the class of measurable
sets \(E\) for which
\[
 x\longmapsto\int\mathbf1_E(\theta,x)\Pi_x(d\theta)
\]
is measurable and its integral against \(\mu\) is \(Q(E)\).
Rectangles belong to \(\mathcal E\). The whole space belongs;
complements preserve both conditions by subtraction from one;
and countable disjoint unions preserve them by monotone convergence.
Thus \(\mathcal E\) is a Dynkin system containing the rectangle
pi-system. The pi-lambda theorem gives all product measurable sets.
Finite sums of their indicators and monotone approximation now prove
\[
 \int h(\theta,x)\,Q(d\theta,dx)
   =\int\left(\int h(\theta,x)\Pi_x(d\theta)\right)\mu(dx)
\]
for every nonnegative product measurable \(h\).
Measurability of the inner integral follows from the same
set argument and simple-function approximation. This verifies the full conditional
kernel property, including data-dependent parameter events.

At \(x\in G\), its defining density proves \(\Pi_x\ll\Pi\)
and the integral formula for indicator functions. Approximation
proves that formula for nonnegative measurable functions.
For signed \(h\) with finite posterior absolute integral, apply
the nonnegative identity to its positive and negative parts and
subtract finite values. No ratio on a zero normalizer has been used.
\end{proof}
\end{annexitem}

\begin{annexitem}{annex:p15.bayes-event-odds}{thm:p15-bayes-event-odds}{All branches of the odds identity}
\begin{proof}
Partitioning the normalizing integral into \(H,H^c\) gives
\(m=p\ell_H+(1-p)\ell_{H^c}\).
On \(G\) both terms are finite, their sum is positive, and the
posterior formula gives the stated probability. When both
likelihoods are positive, division of the two posterior masses yields
\[
 \frac{\Pi_x(H)}{\Pi_x(H^c)}
       =\frac{p\ell_H}{(1-p)\ell_{H^c}}.
\]
The map \(q\mapsto q/(1-q)\) is strictly increasing on \((0,1)\):
cross multiplication reduces its comparison to that of the numerators
\(q(1-r)\) and \(r(1-q)\). The update therefore increases \(p\)
exactly when the multiplier exceeds one.

If only the numerator likelihood is zero, the probability formula
is zero. If only its complement likelihood is zero, it is one.
If both are zero their sum is zero, outside \(G\).
Finally if \(\Pi(H)=0\), absolute continuity of the posterior
on \(G\) gives \(\Pi_x(H)=0\). This last assertion does not
use the conditional-prior expressions, which were defined only
for \(0<\Pi(H)<1\).
\end{proof}
\end{annexitem}

\begin{annexitem}{annex:p15.bayes-flag-calculation}{thm:p15-bayes-flag-calculation}{Base rates, new evidence and repeated information}
\begin{proof}
For one flag, the joint positive masses under \(H,H^c\) are
respectively \(99/100000\) and \(999/100000\). Their normalized
ratio is
\[
              \frac{99}{99+999}=\frac{11}{122}.
\]
The likelihood ratio is \((99/100)/(1/100)=99\).
Under conditional independence, the positive-pair likelihoods
are \(9801/10000\) and \(1/10000\), so
\[
 P(H\mid+,+)=\frac{9801}{9801+999}
            =\frac{9801}{10800}=\frac{363}{400}.
\]
If \(Z_2=Z_1\), the events
\(\{Z_1=+,Z_2=+\}\) and \(\{Z_1=+\}\) coincide. Their
conditional probabilities under either state are the one-flag
probabilities. Thus the posterior remains \(11/122\).
The two constructions can have identical one-flag marginals;
their different joint laws account for the different answers.

For general prior \(p\), normalization gives
\[
       \frac{(99/100)p}{(99/100)p+(1/100)(1-p)}
                          =\frac{99p}{1+98p}.
\]
Its derivative is \(99/(1+98p)^2>0\), proving monotonicity.
At the two endpoints its values are
\(99/10098=1/102\) and \(99/198=1/2\).
All probabilities are positive under the stated assumptions,
so every displayed conditioning operation is defined.
\end{proof}
\end{annexitem}

\begin{annexitem}{annex:p15.bayes-sequential-update}{thm:p15-bayes-sequential-update}{Conditional factors and telescoping evidence}
\begin{proof}
For a fixed finite \(n\), the product of the normalized conditional
densities is a density on the product observation space. Iterated
nonnegative integration, beginning with \(x_n\), gives total integral
one at each parameter. Mixing with \(\Pi\) gives the joint model
and its predictive density \(m_n\). Each prefix marginal is the
same \(m_i\), since integrating the later conditional densities
successively gives one.

At a regular prefix, substitute
\(\Pi_{i-1}(d\theta)=L_{i-1}(\theta)\Pi(d\theta)/m_{i-1}\)
into the predictive integral. This gives
\[
 r_i=\frac{1}{m_{i-1}}\int f_iL_{i-1}\,d\Pi
     =\frac{m_i}{m_{i-1}}.
\]
All terms are finite and positive. Multiplying this density
against \(\Pi_{i-1}\) and dividing by \(r_i\) yields
\(L_i\Pi/m_i=\Pi_i\), proving that successive and batch
updates agree. Multiplication over \(i\) cancels all adjacent
normalizers and leaves \(m_n/m_0=m_n\). Taking logarithms
proves the cumulative identity.

Regularity holds almost surely for the joint observation model
at all prefixes of a fixed finite path: each exceptional prefix
set is null under its marginal and there are only finitely many.
For an infinite sequence, the countable union of these null
prefix events is still null whenever the joint sequence exists.
No claim at an exceptional observation is needed.

If \(f_i(x_i\mid\theta,x_{1:i-1})=f_\theta(x_i)\),
the product is unchanged by a permutation of the data.
The normalizer and posterior are therefore unchanged.
For a mechanism that depends on the preceding observations,
permuting values can change the conditional factors.
The copied-flag example already proves that multiplying
unconditional marginals can produce a different posterior:
the conditional probability of its second positive, given
the first, is one under either state, not the original
one-flag probability.
\end{proof}
\end{annexitem}

\begin{annexitem}{annex:p15.bayes-predictive-kernel}{thm:p15-bayes-predictive-kernel}{The future kernel and every moment component}
\begin{proof}
For fixed \(x\), integrating the probability measure
\(R(\theta,x,\cdot)\) against \(\Pi_x\) preserves nonnegativity,
total mass one and countable additivity. The last property follows
by monotone convergence applied to disjoint event sums.
For fixed future event \(B\), the map
\((\theta,x)\mapsto R(\theta,x,B)\) is measurable. Integration
against the kernel \(x\mapsto\Pi_x\) is measurable by the
simple-function argument in the first annex item.
Thus \(\widehat P\) is a probability kernel.

Construct the joint law of \((\theta,X,Y)\) by adjoining \(R\)
to \(Q\), using \cref{thm:p15-kernel-joint-measure}.
For observation event \(A\) and future event \(B\), its probability is
\[
 \int_{\Theta\times A}R(\theta,x,B)\,Q(d\theta,dx)
 =\int_A\left(\int R(\theta,x,B)\Pi_x(d\theta)\right)\mu(dx).
\]
This is exactly the conditional-law identity.
Indicator approximation gives the corresponding formula for
any nonnegative future function, then for integrable signed
functions by separating their parts.

Fix an \(x\) with finite predictive second moment. The integral
of the conditional second moments against \(\Pi_x\) is finite,
so those moments are finite almost everywhere. Cauchy--Schwarz
gives \(a(\theta,x)^2\leq E(Y^2\mid\theta,x)\).
Thus \(a\) is square integrable under \(\Pi_x\), and all terms
below are finite. Write \(\bar a=\int a\,d\Pi_x\). Using
\(E(Y^2\mid\theta,x)=b(\theta,x)+a(\theta,x)^2\) gives
\[
 E(Y^2\mid x)-\bar a^2
   =\int b\,d\Pi_x+\int a^2\,d\Pi_x-\bar a^2,
\]
which is the variance decomposition.

For a pair, finite second moments imply integrability of its
product by Cauchy--Schwarz. Conditional independence yields
\(E(Y_1Y_2\mid\theta,x)=a_1(\theta,x)a_2(\theta,x)\).
Integration and subtraction of the product of the marginal
predictive means give
\(\operatorname{Cov}(Y_1,Y_2\mid x)
 =\operatorname{Cov}_{\Pi_x}(a_1,a_2)\).
This explains dependence remaining after a shared parameter
has been integrated out.
\end{proof}
\end{annexitem}

\begin{annexitem}{annex:p15.bayes-beta-calculation}{thm:p15-bayes-beta-calculation}{Normalized conjugacy and the complete predictive count}
\begin{proof}
The beta integral and its normalization were proved in
\cref{thm:p15-beta-gamma-convolution}. For a specified word,
the likelihood is \(\theta^k(1-\theta)^{n-k}\).
Multiplying its prior density gives
\[
 \frac{\theta^{\alpha+k-1}(1-\theta)^{\beta+n-k-1}}
      {B(\alpha,\beta)}.
\]
Since \(a,b>0\), its integral is
\(B(a,b)/B(\alpha,\beta)\), finite and positive. Dividing
gives the beta posterior. There are \(\binom nk\) words
with that count, each with the same likelihood. Summing
their joint densities supplies the count factor, which
cancels from the posterior ratio. For \(n=0\), the one
empty word has probability one and the prior is unchanged.

Conditionally on \(\theta\), a future count has probability
\(\binom rj\theta^j(1-\theta)^{r-j}\). Integrating against
the beta posterior yields the stated beta-integral ratio.
Its probabilities sum to one because the binomial theorem
gives
\(\sum_j\binom rj\theta^j(1-\theta)^{r-j}=1\).
The beta recurrence, obtained from the beta-gamma identity,
gives
\[
 E(\theta\mid x)=\frac at,\qquad
 E(\theta^2\mid x)=\frac{a(a+1)}{t(t+1)},\qquad
 \operatorname{Var}(\theta\mid x)=\frac{ab}{t^2(t+1)}.
\]
The conditional binomial mean and variance are \(r\theta\)
and \(r\theta(1-\theta)\); these follow by summing its
independent Bernoulli indicators. Hence
\[
 E(J_r\mid x)=\frac{ra}{t},
\]
and total variance gives
\[
 \begin{split}
 \operatorname{Var}(J_r\mid x)
 &=r\,\frac{ab}{t(t+1)}
       +r^2\,\frac{ab}{t^2(t+1)}\\
 &=\frac{rab(t+r)}{t^2(t+1)}.
 \end{split}
\]
This also proves that using an independent binomial
distribution with parameter \(a/t\) loses the shared-parameter
variance when \(r>1\).

In the worked example, \(a=9,b=3,t=12\).
Its mean is \(9/12=3/4\) and variance
\(27/(144\cdot13)=3/208\).
The logarithmic derivative of its positive density on
\((0,1)\) is \(8/\theta-2/(1-\theta)\), positive below
\(4/5\) and negative above it. Both endpoint limits of the
density are zero, so this is its unique mode.
The two-success probability is
\[
 E(\theta^2\mid x)=\frac{9\cdot10}{12\cdot13}
                             =\frac{15}{26}.
\]
The square of the posterior mean is \(9/16\); their
difference is \(3/208\), exactly the posterior variance.
\end{proof}
\end{annexitem}

\begin{annexitem}{annex:p15.bayes-normal-calculation}{thm:p15-bayes-normal-calculation}{Completing the square with the marginal constants retained}
\begin{proof}
The normal density and its moments are established in
\cref{thm:p15-normal-normalization}.
Since \(\sum_i(x_i-\bar x)=0\), expansion gives
\[
 \sum_i(x_i-\theta)^2=S+n(\bar x-\theta)^2.
\]
The likelihood therefore factors as
\[
 \begin{split}
 \prod_i\phi_{\sigma^2}(x_i-\theta)
 &=h(x_{1:n})\,\phi_{s^2}(\bar x-\theta),\\
 h(x_{1:n})
 &=\frac{(2\pi\sigma^2)^{-(n-1)/2}}{\sqrt n}
                            e^{-S/(2\sigma^2)}.
 \end{split}
\]
Indeed the density in \(\bar x\) has leading factor
\(\sqrt n(2\pi\sigma^2)^{-1/2}\); multiplying by \(h\)
recovers \((2\pi\sigma^2)^{-n/2}\).

The definitions of \(u,v\) imply the exact square identity
\[
 \frac{(\theta-\mu_0)^2}{\tau^2}
  +\frac{(\bar x-\theta)^2}{s^2}
  =\frac{(\theta-u)^2}{v}
       +\frac{(\bar x-\mu_0)^2}{\tau^2+s^2}.
\]
For verification, the coefficients of \(\theta^2\) agree
because \(v^{-1}=\tau^{-2}+s^{-2}\), those of \(\theta\)
agree because \(u/v=\mu_0/\tau^2+\bar x/s^2\), and the
remaining constants agree after putting them over
\(\tau^2s^2(\tau^2+s^2)\).
Also \(v=\tau^2s^2/(\tau^2+s^2)\), so the normalizing
constants satisfy
\[
 \phi_{\tau^2}(\theta-\mu_0)\phi_{s^2}(\bar x-\theta)
  =\phi_{\tau^2+s^2}(\bar x-\mu_0)\phi_v(\theta-u).
\]
Multiplication by \(h\) and integration over \(\theta\)
give the full marginal density stated in the body.
Division leaves \(\phi_v(\theta-u)\), proving the posterior.
Every normalizer is finite and positive for every real data vector.

Given the data, construct an independent posterior parameter
\(\theta\sim\operatorname{normal}(u,v)\) and fresh errors
\(\varepsilon_j\sim\operatorname{normal}(0,\sigma^2)\).
Set \(Y_j=\theta+\varepsilon_j\). This realizes the required
conditional future law. The independent-normal convolution
theorem \cref{thm:p15-normal-sums} gives the one-observation
law and, by averaging the errors, the law of the future average.
For distinct \(i,j\), expansion and independence yield
\(\operatorname{Cov}(Y_i,Y_j\mid x)=\operatorname{Var}(\theta\mid x)=v\).

In the numerical case \(s^2=1\), so \(v=1/2\) and
\(u=(1/2)(0+2)=1\). Adding the fresh error variance gives
\(4+1/2=9/2\). A parameter fixed at \(\bar x=2\) instead
produces mean two and omits posterior parameter variation.
\end{proof}
\end{annexitem}

\begin{annexitem}{annex:p15.bayes-transformation}{thm:p15-bayes-transformation}{Transporting probability and changing coordinates}
\begin{proof}
For any measurable target event \(A\),
\(\{g(\theta)\in A\}=\{\theta\in g^{-1}(A)\}\).
Applying the posterior kernel to that event proves the
pushforward formula. For a bimeasurable bijection \(g\),
let \(\widetilde\Pi=g_\#\Pi\) and
\(\widetilde f_\eta(x)=f_{g^{-1}(\eta)}(x)\).
The pushforward integral identity, proved first for
indicators and then by monotone approximation, gives
\[
 \int\widetilde f_\eta(x)\widetilde\Pi(d\eta)
                       =\int f_\theta(x)\Pi(d\theta)=m(x).
\]
The same identity with a target indicator proves
\(\widetilde\Pi_x=g_\#\Pi_x\). Transporting a future kernel
as \(\widetilde R(\eta,x,B)=R(g^{-1}(\eta),x,B)\)
leaves its posterior integral unchanged as well.

If \(\theta\) is uniform on \((0,1)\), then
\(P(\theta^2\leq y)=\sqrt y\) for \(0<y<1\).
Its derivative is \(1/(2\sqrt y)\), which integrates to
one and has this CDF. If instead \(\eta\) is uniform
and \(\theta=\sqrt\eta\), its CDF is \(P(\theta\leq t)=t^2\),
giving density \(2t\).
After one success, multiplication by \(\theta\) and normalization
give posterior densities \(2\theta\) and \(3\theta^2\).
Their first moments are respectively
\[
 \int_0^1 2\theta^2\,d\theta=\frac23,\qquad
 \int_0^1 3\theta^3\,d\theta=\frac34,
\]
and their second moments are
\[
 \int_0^1 2\theta^3\,d\theta=\frac12,\qquad
 \int_0^1 3\theta^4\,d\theta=\frac35.
\]
The latter are the means of \(\eta\). Using the same word
for the uniform distributions has changed the prior
measure; it is not a failure of reparameterization invariance.
\end{proof}
\end{annexitem}

\begin{annexitem}{annex:p15.bayes-reweighting}{thm:p15-bayes-reweighting}{Posterior weights, sharp bounds and a small-prior counterexample}
\begin{proof}
Direct integration gives
\[
 m^w(x)=\frac{\int wf_\theta(x)\Pi(d\theta)}{\Pi w}
       =m(x)\frac{\Pi_xw}{\Pi w}.
\]
The stated finiteness and positivity conditions therefore
make the new observation regular. Substitution into its
posterior formula cancels the two prior normalizers and
gives \(w\Pi_x/\Pi_xw\).

Let \(A=\int_Hw\,d\Pi_x\) and \(B=\int_{H^c}w\,d\Pi_x\).
Absolute continuity transfers the bounds on \(w\) to
\(\Pi_x\)-almost everywhere, so
\[
 ap_x\leq A\leq bp_x,\qquad
 a(1-p_x)\leq B\leq b(1-p_x).
\]
For nonnegative \(A,B\) with positive sum, \(A/(A+B)\)
increases with \(A\) and decreases with \(B\).
The extremal bounds follow, including \(p_x=0,1\).
The lower bound is attained by \(w=a\) on \(H\), \(w=b\)
on \(H^c\); the reversed choices attain the upper bound.
These weights define proper reweighted priors because
\(0<a\leq\Pi w\leq b<\infty\).

For contamination, linearity gives
\(m^\varepsilon=(1-\varepsilon)m+\varepsilon m_R\).
Its posterior numerator is
\[
 (1-\varepsilon)m(x)\Pi_x(A)+\varepsilon m_R(x)R_x(A)
\]
for each event \(A\). Dividing proves the mixture formula.
To see why a small \(\varepsilon\) is not enough to control
the posterior change, take \(\Theta=\{0,1\}\),
\(\Pi=\delta_0,\ R=\delta_1\), and a binary observation
whose positive probability is \(0<\delta<1\) at parameter zero
and one at parameter one. Given a positive,
\[
       \lambda=\frac{\varepsilon}
                      {(1-\varepsilon)\delta+\varepsilon}.
\]
For fixed \(0<\varepsilon<1\), this tends to one as
\(\delta\downarrow0\). Each positive observation in the
sequence has positive probability under both priors.
The construction proves the lack of a general bound
based only on the size of prior contamination.
\end{proof}
\end{annexitem}

\begin{annexitem}{annex:p15.bayes-coverage}{thm:p15-bayes-coverage}{Which calibration follows from conditioning}
\begin{proof}
The graph indicator \(h(\theta,x)=\mathbf1_{\{\theta\in C(x)\}}\)
is measurable and bounded. Its integral under \(Q\), first
conditioning on the parameter, is
\[
 \int_\Theta\int_{\mathcal X}h(\theta,x)P_\theta(dx)\Pi(d\theta)
                        =\int_\Theta c(\theta)\Pi(d\theta).
\]
Using the posterior kernel instead gives
\[
 \int_{\mathcal X}\int_\Theta h(\theta,x)\Pi_x(d\theta)\mu(dx)
                        =\int_{\mathcal X}\Pi_x(C(x))\mu(dx).
\]
Equality proves the identity. Integrating a lower bound
\(1-\alpha\) gives the average coverage bound since \(\mu\)
is a probability measure. Exact equality of posterior
contents gives equality of the two averages.

For a uniform prior and one Bernoulli observation, the
posterior densities are \(2\theta\) after success and
\(2(1-\theta)\) after failure. Their CDFs on \([0,1]\)
are \(F_1(t)=t^2\) and \(F_0(t)=1-(1-t)^2\).
Solving \(F_x(l)=0.025\) and \(F_x(u)=0.975\)
gives precisely the two displayed intervals. The CDFs
are continuous, so their posterior contents are
\(0.975-0.025=0.95\).

Since \(0.99^2=0.9801>0.975\),
\(\sqrt{0.975}<0.99\). Also
\(\sqrt{0.025}>0.01\), so
\(1-\sqrt{0.025}<0.99\).
Both upper endpoints therefore lie below the fixed
parameter \(0.99\). The coverage is
\[
 (0.99)\mathbf1_{\{0.99\in C(1)\}}
       +(0.01)\mathbf1_{\{0.99\in C(0)\}}=0.
\]
This does not contradict the average identity, which
averages over all parameters with the specified prior.

For the normal posterior in the body, the standardized
variable \((\theta-u)/\sqrt v\) is standard normal.
Its density is positive and continuous, so its CDF
is continuous and strictly increasing from zero to one.
Symmetry gives \(\Phi(-z)=1-\Phi(z)\).
The indicated symmetric interval consequently has mass
\((1-\alpha/2)-\alpha/2=1-\alpha\).
This is a posterior calculation and adds no uniform
sampling-coverage conclusion.
\end{proof}
\end{annexitem}

\begin{annexitem}{annex:p15.bayes-decision-rule}{thm:p15-bayes-decision-rule}{Integrated risk and the stated loss functions}
\begin{proof}
Write the finite action set as \(d_1,\ldots,d_k\).
Each \(\rho_x(d_i)\) is measurable by kernel integration.
The event that \(d_i\) is the first minimizer is specified
by finitely many measurable strict or weak comparisons
between these functions. Thus the tie-resolved minimizer
\(\delta_*(x)\) is measurable.

For any measurable rule \(\delta(x)\), its integrated
expected loss is
\[
 \begin{split}
 R(\delta)&=\int_\Theta\int_{\mathcal X}
                L(\delta(x),\theta)P_\theta(dx)\Pi(d\theta)\\
 &=\int_{\mathcal X}\rho_x(\delta(x))\mu(dx).
 \end{split}
\]
The graph is measurable because the action set is finite,
and the loss is bounded, so both representations are valid.
Pointwise minimization gives \(R(\delta_*)\leq R(\delta)\).
In particular this holds against each constant rule.
Its risk is \(\int L(d,\theta)\Pi(d\theta)\), proving the
expected-value inequality. This compares losses before
observation and does not require every realized observation
to improve an individual decision's outcome.

For the two declarations, put \(p_x=\Pi_x(H)\).
Their risks are \(c_{\rm FP}(1-p_x)\) and
\(c_{\rm FN}p_x\). The first is no greater than the second
exactly when
\(c_{\rm FP}\leq(c_{\rm FP}+c_{\rm FN})p_x\).
Division by the positive sum yields the stated threshold.
At equality both actions minimize risk.
For equal losses the threshold is \(1/2>11/122\);
for losses \(1,20\), it is \(1/21<11/122\), since
\(122<231\). This verifies the worked decision comparison.

For squared-error loss on the real line, finite second
moment implies a finite mean \(\bar\theta\).
For any real action \(d\), expanding about that mean gives
\[
 \int(\theta-d)^2\,\Pi_x(d\theta)
   =\int(\theta-\bar\theta)^2\,\Pi_x(d\theta)
                       +(d-\bar\theta)^2.
\]
The cross term vanishes by the definition of the mean.
The second term has its unique minimum zero at
\(d=\bar\theta\). This is a separate pointwise real-action
result; no unproved measurable-selection theorem for
arbitrary infinite action spaces was used.
\end{proof}
\end{annexitem}

\begin{annexitem}{annex:p15.bayes-log-score}{thm:p15-bayes-log-score}{A complete finite logarithmic-loss comparison}
\begin{proof}
The logarithmic loss of \(p\) itself is finite, since
there are finitely many positive masses.
If \(q(y)=0<p(y)\) for any \(y\), its expected loss is
infinite and the conclusion is immediate.
Otherwise let \(S=\{y:p(y)>0\}\).
For \(t>0\), the function \(t-1-\log t\) has derivative
\(1-1/t\), decreases to zero at \(t=1\), and then
increases. Hence \(\log t\leq t-1\), with equality
only at one. Apply this with \(t=q(y)/p(y)\) and sum:
\[
 \sum_{y\in S}p(y)\log\frac{q(y)}{p(y)}
       \leq\sum_{y\in S}(q(y)-p(y))
       =q(S)-1\leq0.
\]
Negating gives nonnegativity of the loss difference.
If equality holds, every positive-weight logarithm
inequality is an equality and \(q(S)=1\).
Thus \(q(y)=p(y)\) on \(S\) and \(q\) assigns zero
mass off \(S\), which proves \(q=p\).
Conversely that equality plainly gives zero loss difference.
Finally the sequential identity
\(m_n=\prod_i r_i\), on its positive finite domain,
gives \(-\log m_n=-\sum_i\log r_i\).
No physical calibration or identification of the
sampling model follows from this algebraic minimization.
\end{proof}
\end{annexitem}

\begin{annexitem}{annex:p15.bayes-stopping}{thm:p15-bayes-stopping}{Observed-path likelihoods and selected experiments}
\begin{proof}
For a nonnegative measurable likelihood \(L_z\), Bayes's
formula with the displayed record law gives, on any
parameter event \(A\),
\[
 \Pi(A\mid z)=
 \frac{\int_A c(z)L_z(\theta)\Pi(d\theta)}
      {\int c(z)L_z(\theta)\Pi(d\theta)}
 =\frac{\int_A L_z(\theta)\Pi(d\theta)}
      {\int L_z(\theta)\Pi(d\theta)}.
\]
The common factor is finite and strictly positive,
so cancellation is justified. A parameter-dependent
factor stays inside both integrals and cannot in
general be removed. For example, with a uniform
prior on two parameters and an otherwise constant
likelihood, recording probabilities \(1\) and \(1/2\)
give posterior masses \(2/3,1/3\), not \(1/2,1/2\).

A parameter-independent randomized stopping rule on
a finite alphabet gives this structure directly.
If its probability of stopping after the observed
prefix \(x_{1:i}\), given no earlier stop, is
\(a_i(x_{1:i})\), the likelihood of a record stopping
at \(n\) is
\[
 \left\{\prod_{i<n}(1-a_i(x_{1:i}))\right\}
       a_n(x_{1:n})\prod_{i=1}^n
                   f_i(x_i\mid\theta,x_{1:i-1}).
\]
The factors in braces and the terminal stop probability
depend on the observed path but not on the parameter.
Only paths with a positive recording factor are in
the theorem's domain. A possible infinite record must
be treated as a different observation, not assigned
a finite-path formula.

In first-success sampling, the finite record at \(n\)
is exactly the word of \(n-1\) failures and one success.
Its likelihood is \((1-\theta)^{n-1}\theta\).
For every \(\theta\in(0,1)\), the probability of
no success in the first \(n\) trials is \((1-\theta)^n\),
which tends to zero. Thus a finite stopping time occurs
with probability one at every such parameter and
under the beta mixture by bounded convergence.
Multiplication by the beta prior and normalization
give shapes \(\alpha+1,\beta+n-1\).
The same fixed-length word has the identical likelihood.
The fixed-length event of one success is a disjoint
union of \(n\) such words, giving the constant factor \(n\).
This establishes the three posterior identities for
their three declared experiments.

For the selected-only experiment, define its sampling
law at each \(\theta\in(0,1)\) as the original Bernoulli
law conditional on \(X=1\). This conditioning is valid
at every such parameter, and
\[
             P_\theta(X=1\mid X=1)=\theta/\theta=1.
\]
With the prior specified for this experiment, its
likelihood is constant and the posterior equals the prior.
In the experiment consisting of one predesignated
trial, the likelihood of its success is instead
\(\theta\); a uniform prior therefore updates to density
\(2\theta\), the beta law with shapes \(2,1\).
If the observation includes that a predesignated trial
was admitted to the selected group, its admission event
also has likelihood \(\theta\), so that is the latter
inference problem. The distinction is in the observed
experiment and its prior specification, not an
alternative manipulation of the same likelihood.
\end{proof}
\end{annexitem}
